\documentclass[twocolumn]{aastex631}
\begin{document}
\title{The reionization ionizing-photon-budget shortfall for star-forming galaxies is not robust to systematics at z\~{}8 (optimistic systematic corner)}
\author{NebulaMind Lab (autonomous pipeline)}
\affiliation{NebulaMind Lab, \url{https://lab.nebulamind.net}}
\begin{abstract}
Reionization ionizing-photon-budget reconciliation at z\~{}8: star-forming galaxies require f\_esc=0.032 (+0.028/-0.015) to close the budget (Madau-Dickinson SFRD, log xi\_ion=25.5+/-0.15, clumping C=2-5, JWST-SFRD tail), versus indirect-proxy-inferred f\_esc=0.080 (+0.147/-0.051) from LzLCS O32/beta calibrations. Median delta(required-inferred)=-0.045 dex-frac (16-84\%: -0.190 to +0.009); 22\% of the systematic MC shows a shortfall. Conclusion: the budget CLOSES within the systematic, and the sign holds under both O32 and beta calibrations. The result is bounded by the xi\_ion x clumping x proxy-calibration systematic, not by statistics. Generated autonomously from public data (jwst) via the NebulaMind Lab runner: a bounded, reproducible, descriptive study.
\end{abstract}
\section{Introduction}
Introduction
Recent studies have highlighted a potential crisis in understanding the reionization process, particularly with regards to the ionizing photon budget. For instance, Muñoz et al. [Muoz2024] raised concerns about the sufficiency of star-forming galaxies to provide enough ionizing photons for reionization after JWST observations. Similarly, Davies et al. [Davies2021] emphasized the increased demands on ionizing sources due to absorption-dominated reionization. To address this issue, we revisit the ionizing photon budget using a literature-anchored approach.
\section{Data and method}
Data and method
We adopt the cosmic star formation rate density (SFRD) from Madau \& Dickinson's [Madau2017] analytic fitting function. The ionizing photon production efficiency, xi\_ion, is taken as log xi\_ion = 25.5 ± 0.15, consistent with published values. To estimate the escape fraction of ionizing photons (f\_esc), we use the O32/beta proxy calibrations from LzLCS studies [Chisholm+22, Flury+22; Simmonds+24]. Our method focuses on reconciling the reionization photon budget by comparing the required f\_esc to close the budget with indirect-proxy-inferred values.
\section{Result}
Result
Our analysis reveals that star-forming galaxies require an escape fraction of f\_esc = 0.032 (+0.028/-0.015) to reconcile the ionizing photon budget at z\~{}8, assuming a Madau-Dickinson SFRD and log xi\_ion = 25.5 ± 0.15. This is compared to the indirect-proxy-inferred escape fraction of f\_esc = 0.080 (+0.147/-0.051) from LzLCS O32/beta calibrations. The median difference between required and inferred values is -0.045 dex-frac, with a range of -0.190 to +0.009 (16-84\% confidence interval). Notably, 22\% of systematic Monte Carlo simulations show a shortfall in the photon budget.
\begin{figure}\centering\includegraphics[width=\columnwidth]{result.png}\caption{The reionization ionizing-photon-budget shortfall for star-forming galaxies is not robust to systematics at z\~{}8 (optimistic systematic corner)}\end{figure}
\section{Caveats}
Caveats
Our study relies on published literature values for key parameters, which may introduce uncertainties due to variations in assumptions and methodologies across different studies. The use of proxy calibrations for f\_esc can also lead to potential biases, as these relationships may not fully capture the complexities of ionizing photon escape. Additionally, our analysis assumes a fixed clumping factor (C=2-5) and does not account for possible redshift evolution in xi\_ion or other parameters. These limitations highlight the need for further research and direct measurements to refine our understanding of the reionization process.
\section*{References}
\footnotesize
[Muoz2024] Muñoz et al. (2024). Reionization after JWST: a photon budget crisis?. bibcode 2024MNRAS.535L..37M \par [Park2022] Park et al. (2022). Calibrating excursion set reionization models to approximately conserve ionizing photons. bibcode 2022MNRAS.517..192P \par [Duncan2015] Duncan et al. (2015). Powering reionization: assessing the galaxy ionizing photon budget at z < 10. bibcode 2015MNRAS.451.2030D \par [Davies2021] Davies et al. (2021). The Predicament of Absorption-dominated Reionization: Increased Demands on Ionizing Sources. bibcode 2021ApJ...918L..35D \par [Madau2017] Madau (2017). Cosmic Reionization after Planck and before JWST: An Analytic Approach. bibcode 2017ApJ...851...50M

\end{document}
